% ============================================================
% ВВЕДЕНИЕ
% ============================================================
\structure{ВВЕДЕНИЕ}
\addcontentsline{toc}{structure}{Введение}

\subsection*{Актуальность темы}

Развитие методов машинного обучения привело к смещению отрасли от относительно небольших классических моделей к крупным языковым моделям (LLM, Large Language Model), мультимодальным архитектурам и диффузионным генераторам. Число параметров современных моделей выросло с сотен миллионов (BERT~\cite{bert2018}, 2018) до сотен миллиардов (Llama~3 70B~\cite{llamapaper2024}, Falcon~180B~\cite{falcon2023}), а размер хранимых весов --- с сотен мегабайт до сотен гигабайт. Несмотря на тенденцию к повышению плотности возможностей моделей~\cite{capability_density}, state-of-the-art результаты по-прежнему требуют архитектур колоссального масштаба.

Одновременно индустрия уверенно движется в сторону бессерверного инференса (serverless inference) --- модели эксплуатации, при которой облачная платформа автоматически масштабирует сервисы в ответ на изменение нагрузки, включая масштабирование до нуля активных реплик при отсутствии входящих запросов. Данная модель реализована в таких платформах, как Knative~\cite{knative} --- открытый стандарт на базе Kubernetes, поддерживаемый Google, IBM и Red Hat; OpenFaaS --- система развёртывания функций в Kubernetes-кластере; Amazon Web Services Lambda --- управляемая serverless-платформа публичного облака; а также KServe~\cite{kserve} --- стандартизованная платформа для ML-инференса с поддержкой scale-to-zero. Эти платформы позволяют значительно экономить ресурсы кластера: ноды не удерживают вычислительные мощности для простаивающих сервисов, а плата взимается только за фактически обработанные запросы.

Однако serverless-модель вступает в острое противоречие с природой современных ML-моделей. Перед началом обработки запросов контейнер инференса должен загрузить файлы весов из внешнего репозитория или общего хранилища. Для модели весом 15--350~ГБ такая загрузка занимает от нескольких минут до десятков минут даже по высокоскоростному каналу. По оценкам, опубликованным командой BentoML~\cite{bentoml2024cold}, при стандартном подходе загрузка весов LLM-модели размером 15~ГБ может занимать свыше 10 минут, что неприемлемо для большинства реальных сервисов. Таким образом, холодный старт (cold start) --- задержка от момента появления запроса до готовности сервиса его обработать --- становится не техническим неудобством, а фактором, определяющим выполнимость SLA (Service Level Agreement, соглашение об уровне обслуживания) в продуктивных системах.

Актуальность задачи оптимизации холодного старта подтверждается и отраслевыми тенденциями: доля ML-нагрузок в публичных облаках стремительно растёт, а требования к задержке инференса ужесточаются. Пользовательские интерфейсы генеративного ИИ --- чаты, ассистенты, системы рекомендаций --- создают нагрузки с ярко выраженной пиковостью, при которых платформа обязана быстро масштабировать парк реплик. С экономической точки зрения serverless-инференс привлекателен: оператор платит только за реально затраченные вычислительные ресурсы, а не за простаивающие GPU-инстансы. Однако реализация этой экономии требует, чтобы время масштабирования сервиса было значительно меньше времени ожидания пользователя. Для нагрузок реального времени приемлемое время ответа составляет единицы секунд, тогда как задержка холодного старта при загрузке крупных моделей измеряется минутами. Это противоречие делает задачу оптимизации холодного старта одной из центральных в области MLOps и инфраструктуры ML-платформ.

\subsection*{Проблема и противоречие}

Классический подход к загрузке весов в Kubernetes использует init-контейнер \texttt{storage-initializer}, который скачивает файлы из внешнего репозитория (HuggingFace, Ollama, ModelScope) в том \texttt{emptyDir}, после чего основной runtime-контейнер загружает эти веса в GPU-память. При горизонтальном масштабировании каждая новая реплика повторяет загрузку заново, не используя данные, уже скачанные другими репликами или оставшиеся на диске ноды от предыдущих запусков.

Масштаб проблемы нагляден на конкретном примере. Предположим, что burst-событие порождает 10 реплик инференса для модели Llama~3 8B с весами объёмом $\approx$15~ГБ. При использовании \texttt{emptyDir}-подхода каждая реплика независимо скачивает модель из HuggingFace Hub:
\[
  V_{\mathrm{total}} = N \times S_{\mathrm{model}} = 10 \times 15~\text{ГБ} = 150~\text{ГБ}
\]
Это 150~ГБ внешнего трафика за одно burst-событие. При типичной ширине канала 1~Гбит/с (125~МБ/с) каждая реплика тратит $15{,}000 / 125 \approx 120$~с~$\approx$~2~мин на загрузку. Все 10 реплик стартуют с задержкой около 2 минут, конкурируя за один и тот же канал к внешнему репозиторию, что на практике увеличивает реальное время ещё в 2--3 раза. Суммарная нагрузка на внешний канал составляет 150~ГБ --- тогда как при наличии кэша достаточно загрузить модель один раз и распределить чанки внутри кластера.

Существующий паллиативный подход --- NFS-хранилище --- устраняет повторные загрузки из внешнего репозитория: первый под скачивает модель, остальные читают из общей директории. Однако NFS вводит центральную точку отказа (SPOF, Single Point of Failure) и bottleneck пропускной способности: при одновременном чтении десяти репликами через один NFS-сервер сеть и диск сервера становятся узким местом. Деградация или недоступность NFS-сервера делает недоступными все сервисы инференса.

Таким образом, существует противоречие: бессерверное масштабирование требует быстрого старта произвольного числа реплик, тогда как существующие механизмы доставки весов либо порождают избыточный внешний трафик (\texttt{emptyDir}), либо вводят централизованное узкое место (NFS). Требуется механизм, который сохраняет преимущества serverless-масштабирования, но устраняет повторную загрузку уже известных кластеру весов без введения новой единственной точки отказа.

\subsection*{Цель работы}

Целью работы является исследование и разработка архитектуры распределённого кэширования файлов моделей машинного обучения для бессерверного инференса в Kubernetes, позволяющей сократить задержку холодного старта, уменьшить нагрузку на внешние источники данных и повысить доступность системы при горизонтальном масштабировании.

\subsection*{Задачи работы}

Для достижения поставленной цели необходимо решить следующие задачи:

\begin{enumerate}[leftmargin=*]
  \item \textbf{Проанализировать причины холодного старта ML-моделей в Kubernetes} --- исследовать зависимость задержки инициализации от размера весов, описать механизм работы \texttt{storage-initializer} и \texttt{emptyDir}, формализовать декомпозицию времени старта и применить закон Литтла для оценки нагрузки на систему.

  \item \textbf{Исследовать существующие решения} для распределённого хранения и кэширования моделей --- провести сравнительный анализ Alluxio, JuiceFS, Dragonfly, CubeFS и Rook/Ceph по критериям POSIX-совместимости, накладных расходов на метаданные, зависимости от внешней инфраструктуры и применимости к serverless-инференсу.

  \item \textbf{Разработать и оценить промежуточный подход на базе NFS} с механизмами координации параллельных загрузок --- исследовать файловые блокировки (\texttt{flock}), ресурс Lease и алгоритм Leader Election; описать ограничения NFS как централизованного решения.

  \item \textbf{Исследовать механизмы координации в Kubernetes} --- детально проанализировать файловые блокировки (\texttt{flock}) над NFS, ресурс Lease из Coordination API, алгоритм Leader Election из библиотеки \texttt{client-go} и обосновать применимость каждого механизма.

  \item \textbf{Спроектировать архитектуру распределённого chunk-кэша} с локальным хранилищем на каждой worker-ноде --- разработать трёхуровневую схему доступа: local cache hit~$\to$~peer cache hit~$\to$~remote fallback; определить протоколы взаимодействия компонентов.

  \item \textbf{Обосновать выбор Redis} в роли реестра метаданных чанков и активных \texttt{storage-agent} --- исследовать TTL-based heartbeat как механизм discovery, сравнить с gossip-протоколом и mDNS, обосновать конфигурации Redis Sentinel и Redis Cluster.

  \item \textbf{Обосновать применение FUSE} (Filesystem in Userspace) для предоставления runtime-контейнеру стандартного файлового интерфейса --- проанализировать механизм VFS-перехватов, оценить накладные расходы context switch и преимущества прозрачной интеграции с любым ML-framework.

  \item \textbf{Сравнить Unix Domain Socket и TCP} и обосновать их роли в локальной и межнодовой коммуникации --- формализовать разницу в задержке и накладных расходах; обосновать применение UDS для канала \texttt{storage-mounter}~$\to$~\texttt{storage-agent} и TCP для канала между агентами на разных нодах.

  \item \textbf{Реализовать прототип компонентов} \texttt{storage-agent} и \texttt{storage-mounter} --- реализовать локальный chunk-кэш на дисковом хранилище, FUSE-файловую систему с lazy-loading, интерфейс к HuggingFace Hub и частичную интеграцию между компонентами.

  \item \textbf{Сформировать методику оценки эффективности и надёжности} предложенной архитектуры --- построить расчётную модель на базе AMAT, проанализировать сценарии отказов, выявить ограничения и определить направления дальнейшего развития.
\end{enumerate}

\subsection*{Объект и предмет исследования}

\textbf{Объект исследования} --- системы развёртывания и масштабирования ML-инференса в Kubernetes/serverless-средах.

\textbf{Предмет исследования} --- методы кэширования, доставки и виртуализации доступа к файлам весов ML-моделей при холодном старте.

\subsection*{Научная новизна и практическая значимость}

Научная новизна работы состоит в следующем:

\begin{itemize}[leftmargin=*]
  \item Предложена специализированная архитектура, объединяющая FUSE-доступ, локальный chunk-кэш на дисках worker-нод, Redis-реестр метаданных и P2P-доставку чанков для сценария serverless ML-инференса в Kubernetes.
  \item Обосновано предпочтение Redis heartbeat перед gossip-discovery и mDNS для малых и средних кластеров с жёстким требованием к задержке на пути cold start.
  \item Формализовано разделение локальной (UDS) и межнодовой (TCP) коммуникации как инструмент минимизации накладных расходов при гарантии совместимости с Kubernetes networking.
\end{itemize}

\textbf{Практическая значимость}: снижение времени повторного старта модели за счёт попадания в локальный или соседний кэш; уменьшение нагрузки на HuggingFace Hub и внешний канал при масштабировании реплик; возможность использовать локальные диски всех worker-нод как распределённый кэш без выделенного хранилищного кластера; устранение NFS как единственной точки отказа.

\subsection*{Структура работы}

Работа состоит из введения, пяти разделов, заключения, списка литературы и приложений.

В первом разделе проанализирована эволюция размеров ML-моделей, описан механизм холодного старта в Kubernetes и его влияние на SLA.

Второй раздел содержит сравнительный анализ существующих решений для хранения и кэширования моделей и обоснование необходимости специализированной разработки.

В третьем разделе описан промежуточный NFS-этап с алгоритмами координации загрузок, его преимущества и ограничения.

Четвёртый раздел посвящён основной разработке --- P2P/FUSE-архитектуре с описанием компонентов \texttt{storage-agent}, \texttt{storage-mounter} и Redis-реестра.

Пятый раздел содержит методику оценки, расчётную модель задержки, ожидаемые результаты сравнения и анализ ограничений.
